%%%%%%%%%%%%%%%%%%%%%%%%%%%%%%%%%%%%%%%%%%%%%%%%%%%%%%%%%%%%%%%%%%%%%%%%%%%%%%%
%% Abstract --- tradução do resumo, mesma afirmação e mesma força.
%%%%%%%%%%%%%%%%%%%%%%%%%%%%%%%%%%%%%%%%%%%%%%%%%%%%%%%%%%%%%%%%%%%%%%%%%%%%%%%

\begin{resumo}[Abstract]
\begin{otherlanguage*}{english}
Vibration-based predictive maintenance of rotating machinery is a standardised practice,
and the obstacle to its adoption is the cost of instrumentation, which restricts coverage
to critical assets. Low-cost sensor nodes address this obstacle, but the literature
proposing them rarely declares the measurement bandwidth and noise floor that underpin the
reported performance. This work determines what the measurement chain specified for a
device of this class --- a consumer-grade accelerometer sampling at 1~kHz, a low-power
microcontroller, a long-range radio link and a polymer printed enclosure --- is able to
measure, and calibrates the decision stage so that performance is verifiable at the
intended operating point. The method builds reference conditions whose correct result
follows from an analytical solution or from a normative definition, submits them to the
proposed chain and measures the deviation. Results have three provenances, declared
individually: synthesised signal, real test-rig data decimated in software, and modelling
over design specifications. On real data, at a 2\% false alarm rate, detection reaches
75.0\% $\pm$ 3.5\% for imbalance, between 14.8\% and 20.7\% for pure bearing faults and
between 5.0\% and 7.2\% for misalignment. On synthesised signals, three estimation
decisions alter the results by larger margins than the physical effect under
investigation: time-domain integration misses the normative indicator by 73\% without bias
and by 716\% under a bias within the sensor specification, whereas frequency-domain
integration recovers it exactly; the detector's default thresholding policy flags 42.1\% of
normal operation as anomalous, against 1.4\% under a calibrated threshold; and random
window partitioning inflates accuracy by roughly 36 percentage points relative to
partitioning by test run. From modelling, the assembly mounting frequency falls at 557~Hz,
inside the measurement band, and the thermal operating limit is imposed by the battery
rather than by the enclosure. We conclude that the chain supports imbalance detection under
a fixed false alarm rate and does not support bearing fault diagnosis, and that the
verifiable contribution of a device of this class lies in the decision stage, provided the
acquisition stage is declared. The conclusions hold for the specified chain, not for the
performance of the assembled device in industrial operation: building the prototype and
testing it in a real environment constitute the next stage.

\vspace{\onelineskip}
\noindent
\textbf{Keywords}: predictive maintenance; vibration analysis; anomaly detection;
embedded systems; low-cost instrumentation.
\end{otherlanguage*}
\end{resumo}
